% !TEX root = ../main.tex
\chapter{Tate 的论文}\label{ch:06}

\begingroup
\renewcommand{\thefootnote}{\fnsymbol{footnote}}
\begin{center}
Stephen S. Kudla\footnote{本章是 2001 年 3 月在耶路撒冷希伯来大学「自守形式、$L$ 函数与数论」学校上关于 Tate 论文的两次讲演的稍加扩充的版本。}
\end{center}
\endgroup
\setcounter{footnote}{0}

Tate 的论文《数域中的 Fourier 分析与 Hecke 的 zeta 函数》（普林斯顿，1950）最初以第 XV 章的形式印行，收在 Cassels 与 Fr\"ohlich 编的会议论文集《代数数论》（Thompson Book Company，华盛顿，1967）中。Tate 在文中对 Hecke 在两篇论文 [7] 里为其 Gr\"o\ss encharakter 所配的 $L$ 函数，给出了解析延拓与函数方程的优雅而统一的处理。在 Chevalley 的阿代尔／伊代尔框架下使用交换调和分析，其威力远胜 Hecke 所用的经典技巧。\footnote{Tate 在论文中以及 [3, 第 XV 章] 末尾给出了一些有趣的历史评注。}

下一章会勾勒经典自守形式在\emph{非交换}调和分析（即自守表示）下的类似解释。对一般约化群 $G$ 建立自守表示理论，是二十世纪后半叶数学的主要成就之一。当然，这一发展仍在继续，从后面几章 [6] 即可看出。事后看来，Tate 的工作可以视为给出了最简单的连通约化群 $G=\GL(1)$ 的自守表示与 $L$ 函数理论，因此对任何对现代自守表示论感兴趣的人，它仍是基本文献与出发点。

这几篇文章有两个主要目标。第一，循 Tate 的路子，对任意数域上任意 Hecke 特征 $\chi$ 所配的 $L$ 函数 $L(s,\chi)$，统一处理其解析性质、解析延拓与函数方程等。第二，在最简单的情形引入阿代尔群表示论的观点与技巧。例如，从局部结果过渡到整体结果，以及本征分布的唯一性定理所起的作用，在这一行里都相当典型。后面几章只给轮廓，读者应查阅更详尽的文献。

\subsection*{参考文献}

当然，最好的参考文献是 Tate 的论文本身。此外还有许多阐述，包括 Weil [13, 第 VII 章]、Bump [2, 第 3.1 节]、[5]、[9]，以及 Valenza 与 Ramakrishnan 的近著 [11]。在这份压缩的综述里，我大体沿用 Weil 1966 年 Bourbaki 报告《Fonction zeta et distributions》[12, 14] 所勾勒的路子，该报告更强调 Tate 原文中只是隐含的等变分布。Weil 对 [14, 第三卷, p.~448] 的评注也值得一读。

\section{阿代尔、伊代尔与 Dirichlet 特征}

为说明伊代尔与拟特征所扮演的角色，先看经典 Dirichlet 特征的如下描述。

对正整数 $N$，模 $N$ 的经典 Dirichlet 特征 $\underline\chi_N:\mathbb Z\to\mathbb C$ 是这样得到的函数：先取特征 $\chi_N:(\mathbb Z/N\mathbb Z)^\times\to\mathbb C^\times$，把它在 $\mathbb Z/N\mathbb Z$ 上以 $0$ 延拓，再拉回到 $\mathbb Z$。我们把 $\chi_N$ 称为一个 Dirichlet 特征。若 $N\mid M$，则通过投影
\[
\operatorname{pr}_{M,N}:(\mathbb Z/M\mathbb Z)^\times\longrightarrow(\mathbb Z/N\mathbb Z)^\times,\qquad N\mid M
\]
下的拉回，$\chi_N$ 定义了 $(\mathbb Z/M\mathbb Z)^\times$ 的特征 $\chi_M=\chi_N\circ\operatorname{pr}_{M,N}$，从而也有相应的模 $M$ 的经典 Dirichlet 特征 $\underline\chi_M$。系 $\{(\mathbb Z/N\mathbb Z)^\times,\operatorname{pr}_{M,N}\}$ 的逆向极限是紧、全不连通的拓扑群\footnote{回忆一下，按定义
\[
\widehat{\mathbb Z}^\times=\Bigl\{(x_N)\in\prod_N(\mathbb Z/N\mathbb Z)^\times\ \Bigm|\ \operatorname{pr}_{M,N}(x_M)=x_N\Bigr\}.
\]
有限群 $(\mathbb Z/N\mathbb Z)^\times$ 的乘积紧，故闭子群 $\widehat{\mathbb Z}^\times$ 也紧。投影映射 $\operatorname{pr}_N:\widehat{\mathbb Z}^\times\to(\mathbb Z/N\mathbb Z)^\times$ 都连续且满。}
\[
\widehat{\mathbb Z}^\times:=\varprojlim_N(\mathbb Z/N\mathbb Z)^\times,
\]
而每个 Dirichlet 特征 $\chi_N$ 都可视为 $\widehat{\mathbb Z}^\times$ 的一个\emph{连续}特征 $\chi$。$\chi$ 的\emph{导子}是使 $\chi$ 在投影 $\operatorname{pr}_{N_0}:\widehat{\mathbb Z}^\times\to(\mathbb Z/N_0\mathbb Z)^\times$ 的核上平凡的最小 $N_0$。此时 $N_0\mid N$，且 $\chi$ 是某个唯一确定的 Dirichlet 特征 $\chi_{N_0}$ 的拉回。所有 $M\mid N_0$ 的经典 Dirichlet 特征 $\underline\chi_M$ 都对应同一个 $\chi$，其中 $\underline\chi_{N_0}$ 是唯一的\emph{本原}经典 Dirichlet 特征，其余都是\emph{非本原}的。在经典全纯模形式与 $\GL(2)$ 自守表示的字典里，有类似的现象：归一化新形式扮演本原经典 Dirichlet 特征的角色。

回忆：对固定素数 $p$，
\[
\begin{aligned}
\mathbb Z_p&:=\varprojlim_n\mathbb Z/p^n\mathbb Z
=\{\alpha=a_0+a_1p+a_2p^2+\cdots\mid 0\le a_i\le p-1\},\\
\mathbb Z_p^\times&=\{\alpha\in\mathbb Z_p\mid a_0\ne0\},\qquad
\mathbb Q_p:=\Bigl\{\alpha=\sum_i a_ip^i\Bigr\},
\end{aligned}
\]
等等。

回忆：$\mathbb Q$ 的阿代尔环 $\A$ 是关于子环 $\mathbb Z_p$ 的限制乘积\footnote{元素 $(x_p)$ 的几乎全部分量都必须落在子环 $\mathbb Z_p$ 内。}
\[
\A=\mathbb R\times{\prod_p}'\mathbb Q_p,
\]
同时 $\mathbb Q\hookrightarrow\A$，$\alpha\mapsto(\dots,\alpha,\alpha,\alpha,\dots)$ 是对角嵌入，像是离散子环。类似地，伊代尔群是关于子群 $\mathbb Z_p^\times$ 的限制乘积
\[
\A^\times={\prod_v}'\mathbb Q_v^\times,
\]
有对角嵌入 $\mathbb Q^\times\hookrightarrow\A^\times$，并且有分解
\begin{equation*}\tag{1.1}
\A^\times\simeq\mathbb Q^\times\times\mathbb R^\times_+\times\widehat{\mathbb Z}^\times,\qquad x=\alpha tu.
\end{equation*}

由这最后一个分解可见，任何 Dirichlet 特征 $\chi$ 都定义了一个连续特征
\[
\omega:\mathbb Q^\times\backslash\A^\times\longrightarrow\mathbb C^\times,\qquad
\omega(x)=\omega(\alpha tu)=\chi(u).
\]
最一般的这种\emph{拟特征}\footnote{连续同态 $\mathbb C^\times$ 若取值在 $\mathbb C^1$ 中（即为酉的），才称为特征。}形如 $\omega\cdot|\ |^s$，$s\in\mathbb C$；这里 $|x|=|\alpha tu|=t$。于是 Dirichlet 特征恰是有限阶的拟特征。

\begin{annote}
(1.1) 中的 $\simeq$ 是宽松的写法。由乘积公式，$t=|x|$ 唯一确定；但 $(\alpha,u)$ 有 $\pm$ 的二义性：若 $\alpha tu=\alpha'tu'$，则 $\alpha'/\alpha=u/u'\in\mathbb Q^\times\cap\widehat{\mathbb Z}^\times=\{\pm1\}$。因此 $\omega(\alpha tu)=\chi(u)$ 只对\emph{偶}特征良定义；奇特征须再补上 $\operatorname{sgn}$ 这个阿基米德成分，正是下面练习 (i) 的内容。
\par\vspace{0.15em}{\raggedleft 【来源：Cassels--Fr\"ohlich《Algebraic Number Theory》第 II 章；Tate 论文 §1（A）】\par}
\end{annote}

注意，由于
\[
(\mathbb Z/N\mathbb Z)^\times=\prod_p(\mathbb Z/p^n\mathbb Z)^\times,\qquad n=n_p=\operatorname{ord}_p(N),
\]
任何 Dirichlet 特征 $\chi_N$ 都有分解 $\chi_N=\otimes_p(\chi_N)_p$，其中 $(\chi_N)_p$ 是 $(\mathbb Z/p^n\mathbb Z)^\times$ 的特征。这给出相应的 $\widehat{\mathbb Z}^\times$ 的特征 $\chi$ 的\emph{分解}
\[
\chi=\otimes_p\,\chi_p,
\]
其中每个 $\chi_p$ 是与 $(\chi_N)_p$ 相配的 $\widehat{\mathbb Z}_p^\times$ 的特征。注意，对每个位 $v$，即对每个 $\mathbb Q_p$ 与 $\mathbb Q_\infty=\mathbb R$，都有嵌入
\[
\mathbb Q_v^\times\hookrightarrow\mathbb Q^\times\backslash\A^\times,\qquad
x_v\mapsto(\dots,1,1,x_v,1,1,\dots).
\]
于是任何拟特征 $\omega$ 都决定每个 $\mathbb Q_v^\times$ 的拟特征 $\omega_v$，并且有与前一分解相容的分解
\begin{equation*}\tag{1.2}
\omega=\otimes_v\,\omega_v.
\end{equation*}

\emph{练习}：(i) 验证：模 $N$ 的奇经典 Dirichlet 特征 $\underline\chi$（即 $\underline\chi(-a)=-\underline\chi(a)$）给出的特征 $\omega$ 满足 $\omega_\infty(x)=\operatorname{sgn}(x)=x/|x|$。(ii) 设 $\omega$ 与模 $N$ 的经典 Dirichlet 特征 $\underline\chi_N$ 相联系，定出 $\omega_p(p)$，分 $p\nmid N$ 与 $p\mid N$ 两种情形。

\section{拟特征的 $L$ 函数}

对一般数域 $k$（或有限域上的函数域）\footnote{我们排除这一情形，只是为了省去它所需的额外附带说明。}，图像类似。我们采用如下记号：
\[
\begin{aligned}
\cO&=\text{$k$ 中的整数环},\\
\cO_v&=\text{它在非阿基米德位 $v$（$=$ 一个离散赋值环）处的完备化},\\
\mathcal P_v&=\cO_v\ \text{的极大理想},\\
\varpi_v&=\mathcal P_v\ \text{的一个生成元},\\
\mathbb F_v&=\cO_v/\mathcal P_v=v\ \text{处的剩余域},\\
q_v&=|\cO_v/\mathcal P_v|=\text{它的阶},\\
\operatorname{ord}_v&\colon k_v\longrightarrow\mathbb Z\cup\{\infty\},\quad\text{赋值，}\operatorname{ord}_v(\varpi_v)=1,\\
|x|_v&:=q_v^{-\operatorname{ord}_v(x)},\quad\text{归一化绝对值},\\
\cO_v^\times&=\text{单位群},
\end{aligned}
\]
等等；并且 $\cO_v^\times\times\mathbb Z\xrightarrow{\ \sim\ }k_v^\times$，$(u,n)\mapsto u\varpi_v^n$。对阿基米德位 $v$，$k_v\simeq\mathbb R$（相应地 $\mathbb C$），以 $|x|_v$ 记通常的绝对值（相应地，通常绝对值的平方）。

于是有这些限制乘积：阿代尔
\[
\A={\prod_v}'k_v\qquad\text{与伊代尔}\qquad\A^\times={\prod_v}'k_v^\times,
\]
有对角嵌入 $k\hookrightarrow\A$ 与 $k^\times\hookrightarrow\A^\times$。这里的限制乘积指：对 $x=(\dots,x_v,\dots)\in\A$（相应地 $\A^\times$），几乎全部分量 $x_v$ 落在 $\cO_v$（相应地 $\cO_v^\times$）中。有绝对值
\[
|\ |_\A:\A^\times\longrightarrow\mathbb R^\times_+,\qquad |x|=\prod_v|x_v|_v,
\]
它在 $k^\times$ 上平凡（乘积公式）。(1.1) 的分解因 (i) 整体单位与 (ii) 理想类群的非平凡性而变复杂。这些恰是 Hecke 不得不费力克服的困难，而阿代尔形式化则漂亮地处理了它们。无论如何，与前面一样有嵌入
\[
i_v:k_v^\times\hookrightarrow k^\times\backslash\A^\times,\qquad
x_v\mapsto(\dots,1,1,x_v,1,1,\dots).
\]

\medskip\noindent\textbf{定义 2.1.}\quad \emph{拟特征}或 \emph{Hecke 特征}，指连续复特征 $\omega:\A^\times\to\mathbb C^\times$，它在 $k^\times$ 上平凡。

这一\emph{定义}与 Kowalski 文章 [8] 所述经典 Hecke 特征概念的关系，见下面第 5 节。有分解
\[
\omega=\otimes_v\,\omega_v,
\]
其中 $\omega_v=\omega\circ i_v$。

\medskip\noindent\textbf{定义 2.2.}\quad (i) 对非阿基米德位 $v$，拟特征 $\omega_v:k_v^\times\to\mathbb C^\times$ 称为\emph{无分歧}，指它在 $\cO_v^\times$ 上平凡，从而可写成
\[
\omega_v(x)=t_v^{\operatorname{ord}_v(x)}
\]
的形式，其中 $t_v=t_v(\omega_v)\in\mathbb C^\times$。(ii) $k^\times\backslash\A^\times$ 的拟特征 $\omega$ 在 $v$ 处无分歧，指它在 $v$ 处的局部成分 $\omega_v=\omega\circ i_v$ 无分歧。

对任何 $\omega$，存在有限位集 $S=S(\omega)$，含全部阿基米德位\footnote{此处「无分歧」的概念没有定义。}，使得对一切 $v\notin S$，$\omega_v$ 无分歧。例如，在前一节讨论的 $k=\mathbb Q$ 情形，若 $\omega$ 与模 $N$ 的本原经典 Dirichlet 特征 $\underline\chi_N$ 相联系，则 $S(\omega)=\{\infty\}\cup\{p\mid p\mid N\}$。于是与 $\omega$ 相配有复数族 $\{t_v(\omega)\}_{v\notin S}$。例如，对 $s\in\mathbb C$，拟特征 $\omega_s(x)=|x|_\A^s$ 处处无分歧，并决定集合 $\{q_v^{-s}\}$。

\medskip\noindent\textbf{定义 2.3.}\quad 与拟特征 $\omega$ 相配的（部分）$L$ 函数是 Euler 乘积
\begin{equation*}\tag{EP}
L^S(s,\omega)=\prod_{v\notin S}(1-t_v(\omega)q_v^{-s})^{-1},
\end{equation*}
其中 $S\supset S(\omega)$。因子
\[
L_v(s,\omega_v):=(1-t_v(\omega)q_v^{-s})^{-1}
\]
是与各 $\omega_v$ 相配的\emph{局部 $L$ 因子}。

当然，
\[
L^S(s,\omega)=L^S(0,\omega\omega_s),
\]
所以本可在记号中省去 $s$。通常我们宁愿保留 $s$，而改设 $\omega$ 是特征。此时 Euler 乘积在带形 $\operatorname{Re}(s)>1$ 内绝对收敛。

主要目标是证明这些函数的解析性质，即：
\begin{enumerate}[label=\arabic*.,leftmargin=2.2em,itemsep=0.2em]
\item 通过为 $v\in S$ 的素数补上额外的局部 $L$ 因子（例如阿基米德位者），把部分 $L$ 函数 $L^S(s,\omega)$「完备化」；以及
\item 证明完备 $L$ 函数的亚纯解析延拓与函数方程。
\end{enumerate}

这些目标将这样达成：把局部 $L$ 因子（相应地，整体 Euler 乘积）解释为一维复向量空间上两组自然构造的基之间的比例常数——这是分布的向量空间！

\begin{annote}
「局部 $L$ 因子即比例常数」是本篇的枢纽，值得先说清逻辑。对每个位 $v$，空间 $S'(\omega\omega_s)$ 在下面定理 3.4 的意义下是一维的；其中有两个自然向量：一是 zeta 分布 $z(s,\omega)$，二是归一化后的 $z_o(s,\omega)$。比例式 $z=L\cdot z_o$ 就把 $L_v(s,\omega)$ 定义了出来。整体情形完全相同，只是张量积取代了单个空间，(4.3) 是同一句话的全局版本。因此「$L$ 因子」在此不是先验的 Euler 乘积，而是两组基的换算系数；Euler 乘积的地位反而成了待证的性质。
\par\vspace{0.15em}{\raggedleft 【来源：Tate 论文 §2--§3；Weil《Bourbaki 报告》(1966)（A）】\par}
\end{annote}

\section{局部理论}

\subsection*{本征分布}

固定 $k$ 的一个位 $v$，为简化记号写 $F=k_v$。以 $dx$（相应地 $d^\times x$）记 $F$（相应地 $F^\times$）上的 Haar 测度。注意，对 $a\in F^\times$，
\[
d(ax)=|a|\,dx,
\]
其中 $|\cdot|$ 是上面定义的归一化绝对值。注意于是 $d^\times x=\mu|x|^{-1}dx$，$\mu$ 为某正常数。又对 $s\in\mathbb C$，令
\[
\omega_s(x)=|x|^s
\]
为相应的特征。当然，$\omega_0$ 是平凡特征。若 $v$ 非阿基米德，我们写 $\cO=\cO_v$、$\mathcal P=\mathcal P_v$、$q=q_v$ 等。此时特征 $\omega_s$ 无分歧，且只依赖陪集
\[
s+\frac{2\pi i}{\log(q)}\,\mathbb Z.
\]

\begin{annote}
$\omega_s$ 的这个周期正说明「$s$ 可以省去」一语的分量：对非阿基米德 $v$，$\omega_s$ 只在 $\frac{2\pi i}{\log q}\mathbb Z$ 的陪集上变化，故 $L^S(s,\omega)$ 与 $L^S(0,\omega\omega_s)$ 是同一个量。也正因如此，Euler 乘积的收敛性只取决于 $\omega$ 的酉部分与 $\operatorname{Re}(s)$，而与 $s$ 的虚部无关。
\par\vspace{0.15em}{\raggedleft 【来源：Tate 论文 §2；Bump《Automorphic Forms and Representations》3.1（A）】\par}
\end{annote}

令
\[
\begin{aligned}
S(F)&=\text{$F$ 上 Schwartz--Bruhat 函数空间},\\
S(F)'&=\text{$F$ 上缓增分布空间，即连续线性泛函 }\lambda:S(F)\to\mathbb C\text{ 的空间}.
\end{aligned}
\]

\medskip\noindent\textbf{例。}\quad (i) $F$ 非阿基米德时，复值函数 $f$ 属于 $S(F)$，当且仅当存在依赖 $f$ 的整数 $r\ge0$，使
\[
\operatorname{supp}(f)\subset\mathcal P^{-r}\quad\text{（紧支）},\qquad
f=\text{在 }\mathcal P^r\text{ 的陪集上为常数}\quad\text{（局部常值）}.
\]
空间 $S(F)$ 就是所有这类函数组成的复向量空间，而缓增分布是 $S(F)$ 上任意的 $\mathbb C$-线性泛函。(ii) $F=\mathbb R$（相应地 $\mathbb C$）时，$S(\mathbb R)$ 是通常的 Schwartz 空间\footnote{阿基米德理论的一个好参考文献是 Friedlander 的书 [4]；非阿基米德情形见 [1]。}，由连同各阶导数都速降的复值函数组成，例如
\[
f(x)=h(x)e^{-\pi x^2}\qquad\bigl(\text{相应地 }f(x)=h(x,\bar x)e^{-2\pi x\bar x}\bigr),
\]
其中 $h\in\mathbb C[X]$（相应地 $h\in\mathbb C[X,Y]$）为任意多项式。这是 Fr\'echet 空间，缓增分布是它上面的连续线性泛函。

\begin{annote}
非阿基米德情形说「缓增分布就是\emph{任意}线性泛函」，这不是省略。固定 $r$ 后，支在 $\mathcal P^{-r}$ 内、在 $\mathcal P^r$ 的陪集上常值的函数只构成有限维空间；$S(F)$ 是这些有限维空间的正向极限，取最细局部凸拓扑，于是任何线性泛函自动连续。阿基米德情形没有这一性质：$S(\mathbb R)$ 是 Fr\'echet 空间，必须另行要求连续性，才有「缓增分布」这个名字。
\par\vspace{0.15em}{\raggedleft 【来源：Friedlander《Introduction to the Theory of Distributions》；Tate 论文 §2（A）】\par}
\end{annote}

乘法群 $F^\times$ 作用于 $S(F)$：
\[
r(a)f(x)=f(xa),
\]
$x\in F$、$a\in F^\times$；它作用于 $S(F)'$：
\[
\langle r'(a)\lambda,f\rangle=\langle\lambda,r(a^{-1})f\rangle,
\]
其中 $\langle\,,\rangle$ 记 $S(F)'$ 与 $S(F)$ 的配对。

\medskip\noindent\textbf{定义 3.1.}\quad 对 $F^\times$ 的拟特征 $\omega$，令
\[
S'(\omega)=\{\lambda\in S(F)'\mid r'(a)\lambda=\omega(a)\lambda\}
\]
为 $\omega$-\emph{本征分布}空间。

空间 $S'(\omega)$ 可以「几何地」分析。注意 $F^\times$ 作用在加法群 $F=\{0\}\cup F^\times$ 上有两条轨道。与此相配有包含
\[
C_c^\infty(F^\times)\hookrightarrow S(F),
\]
其中 $C_c^\infty(F^\times)$ 是 $F^\times$ 中具紧支集的 $C^\infty$ 函数空间。由对偶，有分布的正合列
\begin{equation*}\tag{3.1}
0\longrightarrow S(F)'_0\longrightarrow S(F)'\longrightarrow C_c^\infty(F^\times)'\longrightarrow 0,
\end{equation*}
其中 $S(F)'_0$ 是支在 $0$ 的分布子空间。此序列与 $F^\times$ 的作用相容。取 $\omega$-本征空间，得序列
\begin{equation*}\tag{3.2}
0\longrightarrow S'(\omega)_0\longrightarrow S'(\omega)\longrightarrow C_c^\infty(F^\times)'(\omega),
\end{equation*}
其中 $S'(\omega)_0$ 是支在 $0$ 的 $\omega$-本征分布空间。首先有如下简单的唯一性结果：

\begin{annote}
(3.1) 的对偶来历值得记住：$C_c^\infty(F^\times)\hookrightarrow S(F)$ 的像在 $S(F)$ 中稠密，故对偶地 $S(F)'\to C_c^\infty(F^\times)'$ 是单射而不是满射，核就是支在 $0$ 的分布。(3.2) 则是把 (3.1) 限制到 $\omega$-本征空间；原书在此没有写末尾的 $\to0$，因为 $C_c^\infty(F^\times)'(\omega)$ 的满射性尚未知——这正是引理 3.2 与定理 3.4 要处理的内容。
\par\vspace{0.15em}{\raggedleft 【来源：Tate 论文 §2；Bump 3.1（A）】\par}
\end{annote}

\medskip\noindent\textbf{引理 3.2.}\quad \emph{空间 $C_c^\infty(F^\times)'(\omega)$ 是一维的，由分布 $\omega(x)d^\times x$ 张成。特别地，对任何 $\lambda\in S'(\omega)$，存在复数 $c$ 使}
\[
\lambda\big|_{C_c^\infty(F^\times)}=c\cdot\omega(x)d^\times x,
\]
\emph{即若 $\operatorname{supp}(f)\subset F^\times$ 紧，则}
\[
\langle\lambda,f\rangle=c\cdot\int_{F^\times}f(x)\omega(x)\,d^\times x.
\]

剩下的问题是定出支在 $0$ 的缓增分布空间 $S(F)'_0$ 与子空间 $S'(\omega)_0$。由
\[
\langle\delta_0,f\rangle=f(0)
\]
定义的 $\delta$ 分布显然 $F^\times$-不变，且支在 $0$。

\medskip\noindent\textbf{引理 3.3.}\quad (i) \emph{若 $F$ 非阿基米德，则}
\[
S(F)'_0=\mathbb C\cdot\delta_0\subset S'(\omega_0),
\]
\emph{其中 $\omega_0$ 是平凡特征。对 $\omega\ne\omega_0$，有 $S'(\omega)_0=0$。}(ii) \emph{若 $F=\mathbb R$，置 $D=\frac{d}{dx}$，则}\footnote{阿基米德情形这里只给出 $F^\times$-有限的分布。}
\[
S(F)'_0=\bigoplus_{k=0}^{\infty}\mathbb C\cdot D^k\delta_0,
\]
\emph{并且}
\[
D^k\delta_0\in S'(\omega),
\]
\emph{其中 $\omega(x)=x^{-k}$。}(iii) \emph{若 $F=\mathbb C$，置 $D=\frac{\partial}{\partial x}$、$\bar D=\frac{\partial}{\partial\bar x}$，则}
\[
S(F)'_0=\bigoplus_{k,l=0}^{\infty}\mathbb C\cdot D^k\bar D^l\delta_0.
\]
\emph{此外}
\[
D^k\bar D^l\delta_0\in S'(\omega),
\]
\emph{其中 $\omega(x)=x^{-k}\bar x^{-l}$。}

基本的局部唯一性结果如下，我们只勾勒其证明。

\medskip\noindent\textbf{定理 3.4.}\quad \emph{对 $F^\times$ 的任何拟特征 $\omega$，}
\[
\dim S'(\omega)=1.
\]

\medskip\noindent\textbf{注。}\quad (i) 当 $S'(\omega)_0=0$，即没有支在 $0$ 的 $\omega$-本征分布时，上面的推理表明 $\dim S'(\omega)\le1$，故只剩存在性。这归结为证明引理 3.2 中的分布 $\omega(x)d^\times x$ 可以延拓成 $S(F)$ 上的 $\omega$-本征分布。这将在下面通过局部 zeta 积分建立。(ii) 当 $S'(\omega)_0\ne0$ 时，问题变成证明分布 $\omega(x)d^\times x$ \emph{不能}延拓成 $\omega$-本征分布。这一更精细的事实同样出自局部 zeta 积分的分析。

\medskip\noindent\textbf{例。}\quad 设 $F$ 非阿基米德。若 $\omega$ 是\emph{分歧}特征，则 $S'(\omega\omega_s)_0=0$，故对一切 $s$，
\[
\dim S'(\omega\omega_s)\le1.
\]
若 $\omega$ 无分歧、$\omega(x)=t^{\operatorname{ord}(x)}$ 如上，则由引理 3.3，
\[
\dim S'(\omega\omega_s)\le
\begin{cases}
2,&\text{若 }t\cdot q^{-s}=1,\\
1,&\text{其它情形}.
\end{cases}
\]

\subsection*{zeta 积分}

对 $F^\times$ 的特征 $\omega$（即酉拟特征），\emph{局部 zeta 积分}
\begin{equation*}\tag{3.3}
z(s,\omega;f):=\int_{F^\times}f(x)\omega\omega_s(x)\,d^\times x
=\int_F f(x)\omega(x)\mu|x|^{s-1}\,dx
\end{equation*}
当 $\operatorname{Re}(s)>0$ 时对一切 $f\in S(F)$ 绝对收敛。这里 $d^\times x=\mu|x|^{-1}dx$ 同上。在此范围内，分布 $f\mapsto z(s,\omega;f)$ 定义了 $S'(\omega\omega_s)$ 的一个非零元素
\[
z(s,\omega)\in S'(\omega\omega_s).
\]

\subsection*{无分歧的局部理论}

设 $F$ 非阿基米德，$\omega$ 无分歧，$\omega(\varpi)=t$。若 $f$ 在 $F^\times$ 中紧支，则积分
\[
\int_{F^\times}f(x)\omega\omega_s(x)\,d^\times x
\]
是整的。于是思路是：用群代数 $\mathbb Z[F^\times]$ 的适当元素杀掉任意 $f\in S(F)$ 的支集；确切地说，用元素
\begin{equation*}\tag{3.4}
\tau=[1]-[\varpi^{-1}],
\end{equation*}
其中 $\varpi$ 是 $\mathcal P$ 的生成元。由于任何 $f\in S(F)$ 在 $0$ 的某个足够小的邻域（比如 $\mathcal P^r$）内为常数，故对 $x\in\mathcal P^{r+1}$ 有
\[
(r(\tau)f)(x)=f(x)-f(x\varpi^{-1})=0.
\]
于是有分布
\[
z_o(s,\omega)\in S'(\omega\omega_s),
\]
定义为
\begin{equation*}\tag{3.5}
\langle z_o(s,\omega),f\rangle:=\int_{F^\times}(r(\tau)f)(x)\,\omega\omega_s(x)\,d^\times x.
\end{equation*}
另一方面，在半平面 $\operatorname{Re}(s)>0$ 内无需杀支集即可收敛，此时
\[
\begin{aligned}
\langle z_o(s,\omega),f\rangle
&=\int_{F^\times}\bigl(f(x)-f(\varpi^{-1}x)\bigr)\omega\omega_s(x)\,d^\times x\\
&=\int_{F^\times}f(x)\omega\omega_s(x)\,d^\times x-\int_{F^\times}f(\varpi^{-1}x)\omega\omega_s(x)\,d^\times x\\
&=\bigl(1-\omega\omega_s(\varpi)\bigr)\cdot z(s,\omega;f)\\
&=(1-tq^{-s})\cdot z(s,\omega;f)\\
&=L(s,\omega)^{-1}\cdot\langle z(s,\omega),f\rangle.
\end{aligned}
\]
严格按分布来说，这表明对 $\operatorname{Re}(s)>0$，
\begin{equation*}\tag{3.6}
z(s,\omega)=L(s,\omega)\cdot z_o(s,\omega),
\end{equation*}
从而右端的表达式给出 $z(s,\omega)$ 到整个 $s$ 平面的亚纯解析延拓。此外，(3.6) 把局部 $L$ 因子解释为一维空间 $S'(\omega\omega_s)$ 的两组自然基之间的比例常数，只要 $L(s,\omega)$ 无极点。

\begin{annote}
$\tau=[1]-[\varpi^{-1}]$ 是 Tate 论证的第一个技巧。$f$ 在 $0$ 附近为常数，故平移一步再相减，支集就被推到 $F^\times$ 里，于是收敛性不再需要 $\operatorname{Re}(s)>0$。代价是 $z_o$ 与 $z$ 相差因子 $1-tq^{-s}$：这正是 $L$ 因子的来源，极点也随之自动从右端移走。换句话说，「延拓」与「公分母」是同一件事的两面。
\par\vspace{0.15em}{\raggedleft 【来源：Tate 论文 §3；Bump 3.1（A）】\par}
\end{annote}

最后，为看出分布 $z_o(s,\omega)$ 何以「自然」，把乘法 Haar 测度 $d^\times x$ 归一化，使单位群 $\cO^\times$ 的体积为 $1$。然后令
\begin{equation*}\tag{3.7}
f^o=\text{$\cO$ 的示性函数},
\end{equation*}
$\cO$ 是 $F$ 的整数环，并计算：
\begin{equation*}\tag{3.8}
\begin{aligned}
\langle z_o(s,\omega),f^o\rangle
&=\int_{F^\times}\bigl(f^o(x)-f^o(\varpi^{-1}x)\bigr)\omega\omega_s(x)\,d^\times x\\
&=\int_{\cO^\times}d^\times x\\
&=1,
\end{aligned}
\end{equation*}
对一切 $s$ 成立。要点是 $r(\tau)f^o$ 是 $\cO^\times$ 的示性函数。于是 $z_o(s,\omega)$ 永不为零，且对一切 $s$ 给出 $S'(\omega\omega_s)$ 的一个基向量。

\medskip\noindent\textbf{注（$L$ 因子作为最大公分母）。}\quad (3.6) 与 (3.8) 合起来表明，$L(s,\omega)$ 恰好给出这族 zeta 积分 (3.3) 的极点。更确切地说，对任何 $f\in S(F)$，函数
\begin{equation*}\tag{3.9}
\frac{z(s,\omega;f)}{L(s,\omega)}=\langle z_o(s,\omega),f\rangle
\end{equation*}
是整的；且对任何给定的 $s$，存在 $f$（具体地 $f^o$）使 (3.9) 在 $s$ 处的值非零。

\medskip\noindent\textbf{例。}\quad 为补完无分歧情形下定理 3.4 的证明，还有一点要核对：不变分布空间 $S'(\omega_0)$ 含 $z_o(0,\omega_0)$ 且维数至多 $2$，实际上却是\emph{一维}的。论证如下。有分布的正合列
\begin{equation*}\tag{3.10}
0\longrightarrow\mathbb C\cdot\delta_0\longrightarrow S(F)'\longrightarrow C_c^\infty(F^\times)'\longrightarrow 0,
\end{equation*}
过渡到 $\omega_0$-本征分布（即 $F^\times$-不变分布），得
\begin{equation*}\tag{3.11}
0\longrightarrow\mathbb C\cdot\delta_0\longrightarrow S'(\omega_0)\longrightarrow\mathbb C\cdot d^\times x.
\end{equation*}
现在的要点是证明：$d^\times x$ 虽在 (3.10) 的商映射的像中，却不在 (3.11) 的商映射的像中；即 $d^\times x$ 在 $S(F)'$ 中的原像\emph{不是}不变的。

具体地说，由
\begin{equation*}\tag{3.12}
\langle\lambda_0,f\rangle:=\langle d^\times x,f-f(0)f^o\rangle
\end{equation*}
定义的分布 $\lambda_0\in S(F)'$ 给出 $d^\times x$ 在 $S(F)'$ 中的一个原像；它显然在 $\cO^\times$ 下不变（因为 $f^o$ 如此），并满足 $\langle\lambda_0,f^o\rangle=0$。分布 $r'(\varpi)\lambda_0$ 与 $\lambda_0$ 在 $C_c^\infty(F^\times)$ 上一致，故相差一个支在 $0$ 的分布，即
\[
r'(\varpi)\lambda_0=\lambda_0+c\delta_0,
\]
$c$ 为某常数。为定出 $c$（若 $\lambda_0$ 真 $F^\times$-不变，$c$ 应为 $0$），我们如上用群代数元素 $\tau$ 在 $f^o$ 上求值：
\begin{equation*}\tag{3.13}
\begin{aligned}
c&=\langle r'(\varpi)\lambda_0-\lambda_0,\,f^o\rangle\\
&=-\langle\lambda_0,\,r(\tau)f^o\rangle\\
&=-\langle d^\times x,\,r(\tau)f^o\rangle\\
&=-\langle z_o(0,\omega_0),\,f^o\rangle\\
&=-1.
\end{aligned}
\end{equation*}
于是，在 $\lambda_0$ 与 $\delta_0$ 张成的 $S(F)'$ 的二维子空间上，$F^\times$ 通过表示
\begin{equation*}\tag{3.14}
\rho(x)=\begin{pmatrix}1&-\operatorname{ord}(x)\\&1\end{pmatrix}
\end{equation*}
作用。无分歧情形下定理 3.4 的证明至此完成。

\begin{annote}
(3.14) 的矩阵是相对有序基 $(\delta_0,\lambda_0)$ 写的：$\delta_0$ 那一列是 $(1,0)$，$\lambda_0$ 那一列是 $(-\operatorname{ord}(x),1)$，故为上三角。若按正文列举的次序 $(\lambda_0,\delta_0)$ 排，矩阵应写成下三角 $\begin{pmatrix}1&0\\-\operatorname{ord}(x)&1\end{pmatrix}$。这个细节无伤大雅，但抄录时极易写反。
\par\vspace{0.15em}{\raggedleft 【来源：原书 p.~120（PDF 125）算式 (3.14)（A）】\par}
\end{annote}

\subsection*{分歧的局部理论}

设 $F$ 非阿基米德，$\omega$ 分歧，从而在 $\cO^\times$ 上非平凡。$\omega$ 的\emph{导子} $c(\omega)$ 是最小的整数 $c$，使 $\omega$ 在 $1+\mathcal P^c$ 上平凡。由于任何 $f\in S(F)$ 在 $0$ 的邻域内为常数，积分
\[
\int_{F^\times-\mathcal P^n}f(x)\omega\omega_s(x)\,d^\times x
\]
当 $n$ 充分大时与 $n$ 无关。这给出 $z(s,\omega;f)$ 到整个 $s$ 平面的\emph{整}解析延拓，从而对一切 $s$ 给出一维空间 $S'(\omega\omega_s)$ 的基向量
\begin{equation*}\tag{3.15}
z_o(s,\omega):=z(s,\omega).
\end{equation*}
分歧情形下定理 3.4 的证明至此完成。此情形我们置
\begin{equation*}\tag{3.16}
L_v(s,\omega)=1,
\end{equation*}
并注意，对
\begin{equation*}\tag{3.17}
f^o(x)=
\begin{cases}
\omega(x)^{-1},&x\in\cO^\times,\\
0,&\text{其它},
\end{cases}
\end{equation*}
有
\[
\langle z_o(s,\omega),f^o\rangle=1.
\]
于是 gcd 性质 (3.9) 仍成立。

\subsection*{阿基米德情形}

此时造成 zeta 积分 $z(s,\omega;f)$ 极点的是 $f$ 在 $0$ 的整条 Taylor 级数，而非非阿基米德情形中出现的单值 $f(0)$；参看 [2, 命题 3.1.7, p.~271]。用分部积分容易得到 $z(s,\omega;f)$ 的解析延拓。这里只概述结果，采用 [13] 的归一化。为得到特别简洁的表达式，放弃 $\omega$ 酉的约定。

若 $F=\mathbb R$，任何拟特征都形如 $\omega\omega_s$，$s\in\mathbb C$，且
\begin{equation*}\tag{3.18}
\omega(x)=x^{-a},\qquad a=0,\ 1.
\end{equation*}
令
\begin{equation*}\tag{3.19}
L(s,\omega)=\pi^{-\frac{s}{2}}\,\Gamma\Bigl(\frac{s}{2}\Bigr).
\end{equation*}
若 $F=\mathbb C$，任何拟特征都形如 $\omega\omega_s$，$s\in\mathbb C$，且
\begin{equation*}\tag{3.20}
\omega(x)=x^{-a}\bar x^{-b},\qquad a,b\in\mathbb Z,\quad \min(a,b)=0.
\end{equation*}
令
\begin{equation*}\tag{3.21}
L(s,\omega)=(2\pi)^{1-s}\,\Gamma(s).
\end{equation*}

\begin{annote}
(3.19) 与 (3.21) 看似与 $a$、$b$ 无关，其实是「已知 $\omega$ 形如 $x^{-a}$」这一约定下的结果：$\omega\omega_s=|\ |^{s-a}\operatorname{sgn}^a$，按标准公式 $\pi^{-\frac{u+a}{2}}\Gamma\bigl(\frac{u+a}{2}\bigr)$ 取 $u=s-a$，恰好得 $\pi^{-\frac{s}{2}}\Gamma\bigl(\frac{s}{2}\bigr)$。但算 $L(1-s,\omega^{-1})$ 时须用一般公式，不可直接代 (3.19)。同理 $\mathbb C$ 情形须记住 $|z|_\mathbb C=z\bar z$、$dx$ 是 Lebesgue 测度的两倍，否则 (3.21) 会差因子 $2\pi$。
\par\vspace{0.15em}{\raggedleft 【来源：Tate 论文 §4--§5；Weil《Basic Number Theory》第 VII 章（A）】\par}
\end{annote}

\medskip\noindent\textbf{命题 3.5.}\quad (i) \emph{分布}
\[
z_o(s,\omega):=L(s,\omega)^{-1}z(s,\omega)
\]
\emph{有到整个 $s$ 平面的整解析延拓，且对一切 $s$ 给出空间 $S'(\omega\omega_s)$ 的基向量。}(ii) \emph{若}
\[
f^o(x)=
\begin{cases}
f_a(x):=x^ae^{-\pi x^2},&\text{若 }F=\mathbb R,\\
f_{a,b}(x):=x^a\bar x^be^{-2\pi x\bar x},&\text{若 }F=\mathbb C,
\end{cases}
\]
\emph{则}
\[
\langle z_o(s,\omega),f^o\rangle=1.
\]

注意我们用的测度是 $d^\times x=|x|^{-1}dx$，其中 $|x|$ 是通常绝对值（相应地 $|x|=x\bar x$），$dx$ 在 $F=\mathbb R$（相应地 $\mathbb C$）时是 Lebesgue 测度（相应地 Lebesgue 测度的两倍）。(ii) 见 [13, 引理 8, p.~127]。

\medskip\noindent\textbf{注。}\quad 实际上，zeta 分布 $z(s,\omega)$ 在 $s$ 中有亚纯解析延拓，在若干非正整数处有单极点。例如 $F=\mathbb R$ 时，极点在点 $-r\in\mathbb Z_{\le0}$，其中 $r\equiv a\pmod 2$，留数为
\begin{equation*}\tag{3.22}
\operatorname*{Res}_{s=-r}z(s,\omega)=c_rD^r\delta_0,
\end{equation*}
$c_r$ 为非零常数，$D$ 如引理 3.3。$z(s,\omega)$ 在 $s=-r$ 处 Laurent 展开的常数项，给出分布 $x^{-r}d^\times x$ 在 $C_c^\infty(F^\times)$ 上到 $S(F)$ 的延拓。像 (3.9) 那样的简单计算表明，此延拓\emph{不}属于空间 $S'(\omega\omega_{-r})$。$F=\mathbb C$ 情形，$s=-r$ 处的留数落在 $S'(\omega\omega_{-r})$ 中，其中 $\omega\omega_{-r}(x)=x^{-a-r}\bar x^{-b-r}$，且是 $D^{a+r}\bar D^{b+r}\delta_0$ 的非零倍。

\subsection*{Fourier 变换}

固定加法群 $F=F^+$ 的一个非平凡特征 $\psi$，并通过
\[
F\xrightarrow{\ \sim\ }\widehat F,\qquad y\mapsto\bigl(x\mapsto\psi(xy)\bigr)
\]
把 $F$ 与其拓扑对偶特征群
\[
\widehat F:=\operatorname{Hom}_{\mathrm{cont}}(F,\mathbb C^1)
\]
等同。$F$ 非阿基米德时，$\psi$ 的\emph{导子} $\nu(\psi)$ 是使 $\psi$ 在 $\mathcal P^{-\nu}$ 上平凡的最大整数 $\nu$。

函数 $f\in S(F)$ 的 Fourier 变换
\[
\hat f(x)=\int_F f(y)\,\psi(xy)\,dy
\]
有定义，且仍属于 $S(F)$。（这是 Schwartz--Bruhat 函数空间的基本性质。）映射 $f\mapsto\hat f$ 是同构 $\widehat{\ }:S(F)\xrightarrow{\sim}S(F)$。Haar 测度有唯一选择，即关于 $\psi$ 的\emph{自对偶测度}，使 Fourier 反演给出
\[
\widehat{\hat f}(x)=f(-x).
\]
设 $\psi$ 已给定，从此固定 $dx$ 的这一选择。分布的 Fourier 变换定义为
\[
\langle\hat\lambda,f\rangle=\langle\lambda,\hat f\rangle.
\]

\medskip\noindent\textbf{引理 3.6.}\quad \emph{若 $\lambda\in S'(\omega)$ 是 $\omega$-本征分布，则}
\[
\hat\lambda\in S'(\omega^{-1}\omega_1),
\]
\emph{即 $\hat\lambda$ 是 $\omega^{-1}\omega_1$-本征分布，其中 $\omega_1(x)=|x|$。}

\emph{证明。}\quad 留作练习！\hfill$\square$

由唯一性结果（定理 3.4）推知，$z_o(s,\omega)$ 是 $z_o(1-s,\omega^{-1})$ 的常数倍（常数依赖 $s$ 与 $\psi$）。

\medskip\noindent\textbf{推论 3.7（局部函数方程）。}\quad
\begin{equation*}\tag{3.23}
\widehat{z_o(1-s,\omega^{-1})}=\epsilon(s,\omega,\psi)\,z_o(s,\omega),
\end{equation*}
\emph{其中 $\epsilon(s,\omega,\psi)$ 为非零常数，即局部 epsilon 因子，依赖 $s$、$\omega$ 与 $\psi$。}

\begin{annote}
(3.23) 的帽号覆盖整个 $z_o(1-s,\omega^{-1})$，表示分布的 Fourier 变换；若只把帽号盖在 $z_o$ 上，就会读成 \(z_o\) 本身的变换而丢掉 $1-s$ 与 $\omega^{-1}$ 的替换。另须注意 $\widehat{z_o(1-s,\omega^{-1})}$ 与 $z_o(s,\omega)$ 都是 $S'(\omega\omega_s)$ 中的向量：帽子把 $(1-s,\omega^{-1})$ 搬回 $(s,\omega)$，这正是函数方程的分布形式。
\par\vspace{0.15em}{\raggedleft 【来源：Tate 论文 §4；Weil《Bourbaki 报告》(1966)（A）】\par}
\end{annote}

\medskip\noindent\textbf{注。}\quad 对给定的 $f\in S(F)$，局部函数方程可写成关系
\begin{equation*}\tag{3.24}
\frac{z(1-s,\omega^{-1};\hat f)}{L(1-s,\omega^{-1})}
=\epsilon(s,\omega,\psi)\cdot\frac{z(s,\omega;f)}{L(s,\omega)},
\end{equation*}
或等价地写成
\begin{equation*}\tag{3.25}
z(1-s,\omega^{-1};\hat f)
=\epsilon(s,\omega,\psi)\,\frac{L(1-s,\omega^{-1})}{L(s,\omega)}\cdot z(s,\omega;f),
\end{equation*}
其中比例常数
\begin{equation*}\tag{3.26}
\gamma(s,\omega,\psi)=\epsilon(s,\omega,\psi)\,\frac{L(1-s,\omega^{-1})}{L(s,\omega)}
\end{equation*}
现在可以带零点与极点，称为\emph{局部伽马因子}。这是更传统的写法，而推论 3.7 的写法在 [12] 中受强调。

$\epsilon$ 因子 $\epsilon(s,\omega,\psi)$ 的若干重要性质可由关系 (3.23) 得到。例如，若 $f^o$ 是关于 $\omega$ 的标准函数，则在 $f^o$ 上对 (3.23) 两端求值，得有用恒等式
\begin{equation*}\tag{3.27}
\epsilon(s,\omega,\psi)=\langle z_o(1-s,\omega^{-1}),\widehat{f^o}\rangle.
\end{equation*}
按构造，zeta 分布满足
\[
z(s,\omega\omega_t)=z(s+t,\omega).
\]
局部 $L$ 函数也须有同一性质，故同样有
\[
z_o(s,\omega\omega_t)=z_o(s+t,\omega).
\]
标准函数 $f^o$ 在乘 $\omega_t$ 下不变。这意味着
\begin{equation*}\tag{3.28}
\epsilon(s,\omega\omega_t,\psi)=\epsilon(s+t,\omega,\psi).
\end{equation*}
对 $\beta\in F^\times$，令 $\psi_\beta(x)=\psi(\beta x)$。若 $dx$ 是关于 $\psi$ 的自对偶测度，则 $|\beta|^{\frac12}dx$ 是关于 $\psi_\beta$ 的自对偶测度。若 $\hat f$ 是 $f\in S(F)$ 关于 $\psi$ 的 Fourier 变换，则 $|\beta|^{\frac12}r(\beta)\hat f$ 是关于 $\psi_\beta$ 的 Fourier 变换，这对自对偶测度总是如此。于是把它代入 (3.27) 的右端，得
\begin{equation*}\tag{3.29}
\epsilon(s,\omega,\psi_\beta)=|\beta|^{s-\frac12}\omega(\beta)\,\epsilon(s,\omega,\psi).
\end{equation*}

下面给出 $\epsilon$ 因子的显式值。利用关系 (3.28) 与 (3.29)，只需在 $\psi$ 与 $\omega$ 都取「标准」形式时计算这些因子。

回忆：非阿基米德情形，$\nu=\nu(\psi)$ 是 $\psi$ 的导子。阿基米德情形，「标准」选择是
\begin{equation*}\tag{3.30}
\psi(x)=
\begin{cases}
e(x),&\text{若 }F=\mathbb R,\\
e(x+\bar x),&\text{若 }F=\mathbb C,
\end{cases}
\end{equation*}
其中对实数 $t$，$e(t)=e^{2\pi it}$；其自对偶测度就是命题 3.5 之后描述的那些测度。

\medskip\noindent\textbf{命题 3.8.}\quad (i) \emph{设 $F$ 非阿基米德、$\omega$ 无分歧、$\nu(\psi)=\nu$，则}
\[
\epsilon(s,\omega,\psi)=\omega(\varpi^\nu)\,q^{(\frac12-s)\nu}.
\]
(ii) \emph{设 $F$ 非阿基米德、$\omega$ 分歧、导子为 $c$、$\nu(\psi)=\nu$，则}
\[
\epsilon(s,\omega,\psi)=\omega(\varpi^{\nu+c})\,q^{(\frac12-s)(\nu+c)}\,\mathfrak g(\omega,\psi),
\]
\emph{其中 $\mathfrak g(\omega,\psi)$ 是高斯和：}
\[
\mathfrak g(\omega,\psi)=q^{\frac12(\nu+c)}\int_{\cO^\times}\omega^{-1}(y)\,\psi(\varpi^{-\nu-c}y)\,dy,
\]
\emph{$dy$ 是关于 $\psi$ 的 $F$ 上自对偶测度。}(iii) \emph{设 $F=\mathbb R$，写 $\omega(x)=x^{-a}$，$a=0$ 或 $1$，$\psi$ 如 (3.30)。则}
\[
\epsilon(s,\omega,\psi)=i^a.
\]
(iv) \emph{设 $F=\mathbb C$，写 $\omega(x)=x^{-a}\bar x^{-b}$，$a,b\in\mathbb Z$、$\min(a,b)=0$，$\psi$ 如 (3.30)。则}
\[
\epsilon(s,\omega,\psi)=i^{\max(a,b)}.
\]

\emph{证明。}\quad 由 (3.27)，须计算关于 $\omega$ 的标准函数 $f^o$ 的 Fourier 变换，其中特征 $\psi$ 可取为标准形。例如情形 (i) 中 $f^o$ 是 $\cO$ 的示性函数，设 $\nu(\psi)=0$，则 $\widehat{f^o}=f^o$。此情形给出 $\epsilon(s,\omega,\psi)=1$。一般无分歧情形随即由 (3.29) 得到。

情形 (ii) 中 $f^o$ 由 (3.17) 给出。于是，依照 [13, 命题 13, p.~131]，
\[
\widehat{f^o}(x)=\int_{\cO^\times}\omega^{-1}(y)\,\psi(xy)\,dy
\]
除非 $\operatorname{ord}(x)=-\nu-c$（$c=c(\omega)$ 为 $\omega$ 的导子）否则为零。事实上
\begin{equation*}\tag{3.31}
\widehat{f^o}(x)=\mathfrak g(\omega,\psi)\,q^{-\frac12(\nu+c)}\,\overline{f^o(\varpi^{\nu+c}x)},
\end{equation*}
其中
\begin{equation*}\tag{3.32}
\mathfrak g(\omega,\psi)=q^{\frac12(\nu+c)}\int_{\cO^\times}\omega^{-1}(y)\,\psi(\varpi^{-\nu-c}y)\,dy.
\end{equation*}
把 (3.32) 用两次得 $\mathfrak g(\omega,\psi)\overline{\mathfrak g(\omega,\psi)}=1$。由于 $\widehat{f^o}$ 是关于 $\omega^{-1}$ 的标准函数，得
\[
\epsilon(s,\omega,\psi)=\omega(\varpi^{\nu+c})\,\mathfrak g(\omega,\psi)\,q^{(\frac12-s)(\nu+c)}.
\]

情形 (iii) 中，对 $f^o=f_a$（如命题 3.5）计算：
\begin{equation*}\tag{3.33}
\begin{aligned}
\widehat{f_a}(x)&=\int_{\mathbb R}y^a\,e^{-\pi y^2}\,e(xy)\,dy\\
&=(2\pi i)^{-a}D_x^a\Bigl\{\int_{\mathbb R}e^{-\pi y^2}\,e(xy)\,dy\Bigr\}\\
&=(2\pi i)^{-a}D_x^a\{e^{-\pi x^2}\}\\
&=i^a\,f_a(x).
\end{aligned}
\end{equation*}
由于 $f_a$ 也是关于 $\omega^{-1}=\omega\omega_2$ 的标准函数，得 (iii)。

\begin{annote}
「$\omega^{-1}=\omega\omega_2$」只在 $a=1$ 时字面成立（此时 $x^{-1}|x|^2=x$）；$a=0$ 时 $\omega^{-1}=\omega=1$。正确的等式是 $\omega^{-1}=\omega\omega_{2a}$。真正的依据是：$f_a=x^ae^{-\pi x^2}$ 是关于拟特征 $\omega^{-1}\omega_{1-s}=|\ |^{1-s+a}\operatorname{sgn}^a$ 的标准函数，故 $z(1-s,\omega^{-1};f_a)=L(1-s,\omega^{-1})$；与 $\widehat{f_a}=i^af_a$ 合用即得 $\epsilon=i^a$。原文此处漏了一个 $a$。
\par\vspace{0.15em}{\raggedleft 【来源：原书 p.~125（PDF 130）证明 (iii)；Tate 论文 §4（A）】\par}
\end{annote}

情形 (iv) 中 $f^o=f_{a,b}$，如命题 3.5。于是，注意 $\min(a,b)=0$，
\begin{equation*}\tag{3.34}
\begin{aligned}
\widehat{f_{a,b}}(x)&=\int_{\mathbb C}y^a\bar y^b\,e^{-2\pi y\bar y}\,e(xy+x\bar y)\,dy\\
&=(2\pi i)^{-a}D_x^a\,(2\pi i)^{-b}\bar D_x^b
\Bigl\{\int_{\mathbb C}e^{-2\pi y\bar y}\,e(xy+x\bar y)\,dy\Bigr\}\\
&=(2\pi i)^{-a}D_x^a\,(2\pi i)^{-b}\bar D_x^b\{e^{-2\pi x\bar x}\}\\
&=i^{a+b}\,f_{b,a}(x).
\end{aligned}
\end{equation*}
注意 $f_{b,a}$ 是关于 $\omega^{-1}(x)=\bar x^{-a}x^{-b}|x|^{a+b}$ 的标准函数，即得所断言的值。更多细节见 Tate 论文或 [13]。\hfill$\square$

\section{整体理论}

\subsection*{分布}

回到带阿代尔环 $\A$ 的数域 $k$，我们定义 $\A$ 上 Schwartz--Bruhat 函数空间 $S(\A)$。此空间含所有形如
\[
f=\otimes_v\,f_v
\]
的函数，其中对一切 $v$ 有 $f_v\in S(k_v)$，且对几乎所有的 $v$ 有 $f_v=f_v^o$，即 $\cO_v$ 的示性函数。注意
\[
f(x)=\prod_v f_v(x_v)
\]
于是有定义，因为几乎所有因子都是 $1$。这类可分解函数的有限线性组合在 $S(\A)$ 中稠密。又 $S(\A)'$ 是缓增分布空间，即 $S(\A)$ 上的连续线性泛函。

\medskip\noindent\textbf{引理 4.1.}\quad \emph{设 $\{\lambda_v\}$ 是一族局部缓增分布 $\lambda_v\in S(k_v)'$，且对几乎所有的 $v$ 有 $\langle\lambda_v,f_v^o\rangle=1$。则存在唯一的 $\A$ 上缓增分布 $\lambda=\otimes_v\lambda_v$，它在可分解函数上由}
\[
\langle\lambda,f\rangle=\prod_v\langle\lambda_v,f_v\rangle
\]
\emph{给出。}

反之，若 $\lambda\in S(\A)'$ 非零，选可分解的 $f$ 使 $\lambda(f)=1$。则对任何位 $v$，写
\[
f=f_v\otimes f^v,\qquad f^v=\otimes_{w\ne v}f_w,
\]
并定义映射
\[
S(k_v)\longrightarrow S(\A),\qquad g_v\mapsto g_v\otimes f^v,
\]
即只变动第 $v$ 个分量。于是得到局部分布
\[
\langle\lambda_v,g_v\rangle:=\langle\lambda,g_v\otimes f^v\rangle,
\]
且对几乎所有的 $v$ 有 $\langle\lambda_v,f_v^o\rangle=1$，并有 $\lambda=\otimes_v\lambda_v$。

伊代尔群 $\A^\times$ 作用于 $S(\A)$ 与 $S(\A)'$；对 $\A^\times$ 上在 $k^\times$ 上取平凡值的拟特征 $\omega$，以 $S'(\omega)$ 记 $\omega$-本征分布空间。整体拟特征的分解 $\omega=\otimes_v\omega_v$ 给出分解
\[
S'(\omega)=\otimes_v\,S'_v(\omega_v),
\]
其中我们写 $S'_v=S(k_v)'$。局部唯一性定理 3.4 随即给出整体唯一性结果：

\medskip\noindent\textbf{定理 4.2.}\quad \emph{对 $\A^\times$ 的任何拟特征 $\omega$，整体 $\omega$-本征分布空间 $S'(\omega)$ 的维数为 $1$。对任何 $s\in\mathbb C$，空间 $S'(\omega\omega_s)$ 由标准的 $\omega\omega_s$-本征分布}
\[
z_o(s,\omega):=\otimes_v\,z_o(s,\omega_v)
\]
\emph{张成。}

\subsection*{整体 zeta 积分与完备 $L$ 函数}

对特征 $\omega$（即 $\A^\times$ 上的酉拟特征）与函数 $f\in S(\A)$，可以定义整体 zeta 积分
\begin{equation*}\tag{4.1}
z(s,\omega;f)=\int_{\A^\times}f(x)\,\omega\omega_s(x)\,d^\times x.
\end{equation*}
若 $f$ 可分解，它是相应局部积分之积：
\begin{equation*}\tag{4.2}
z(s,\omega;f)=\prod_v z(s,\omega_v;f_v).
\end{equation*}
这里的关键是：此积分\emph{收敛}，恰在局部积分之积收敛时。由于存在有限位集 $S$，使 $v\notin S$ 时 $\omega_v$ 无分歧且 $f_v=f_v^o$，故只要乘积
\[
L^S(s,\omega)=\prod_{v\notin S}L_v(s,\omega_v)
\]
绝对收敛，即 $\operatorname{Re}(s)>1$，就得到收敛性。

按分布来说，在半平面 $\operatorname{Re}(s)>1$ 内有 $z(s,\omega)=\otimes_vz(s,\omega_v)$，并且
\begin{equation*}\tag{4.3}
z(s,\omega)=\Lambda(s,\omega)\cdot z_o(s,\omega),
\end{equation*}
其中我们定义\emph{完备 $L$ 函数}
\begin{equation*}\tag{4.4}
\Lambda(s,\omega):=\prod_v L_v(s,\omega),
\end{equation*}
此时包含全部局部 $L$ 因子；而
\begin{equation*}\tag{4.5}
z_o(s,\omega)=\otimes_v z_o(s,\omega_v)
\end{equation*}
就是上面定义的标准 $\omega\omega_s$-本征分布。于是在收敛半平面内，完备 $L$ 函数可以看作两个「自然」整体本征分布之间的比例因子。

\subsection*{Fourier 变换与整体函数方程}

固定 $\A$ 的一个在 $k$ 上平凡的整体加法特征 $\psi$。限制到局部成分，可以写
\[
\psi(x)=\prod_v\psi_v(x_v).
\]
对每个非阿基米德位 $v$，$\psi_v$ 的导子 $\nu_v:=\nu(\psi_v)$ 有定义。于是得到 $\A$ 与其连续特征群的一个等同，它与局部情形中各 $\psi_v$ 所定义的等同相容。我们定义整体 Fourier 变换
\[
\widehat{\ }:S(\A)\xrightarrow{\ \sim\ }S(\A),
\]
它与局部者相容，即若 $f$ 可分解，则
\[
\hat f=\otimes_v\hat f_v.
\]
$\A$ 上关于 $\psi$ 自对偶的测度 $dx$ 可以写成各 $k_v$ 上关于 $\psi_v$ 自对偶的测度之积。分布的 Fourier 变换如前定义。

完备 $L$ 函数 $\Lambda(s,\omega)$ 所需的解析延拓与函数方程，将由以下三者的结合得到：(i) 比例关系 (4.3)，(ii) 各 $z_o(s,\omega_v)$ 的局部函数方程，以及 (iii) 下面的整体恒等式：

\medskip\noindent\textbf{定理 4.3.}\quad \emph{由整体 zeta 积分 (4.1) 在 $\operatorname{Re}(s)>1$ 中定义的分布 $z(s,\omega)$，有到整个 $s$ 平面的亚纯解析延拓，并满足函数方程}
\[
\widehat{z(1-s,\omega^{-1})}=z(s,\omega).
\]

假设定理 4.3，由 (4.3) 立即得到 $\Lambda(s,\omega)$ 的相应解析延拓，因为「归一化」的整体分布 $z_o(s,\omega)$ 是整的。回忆：局部函数方程形如 (3.13)：
\[
\widehat{z_o(1-s,\omega_v^{-1})}=\epsilon_v(s,\omega_v,\psi_v)\,z_o(s,\omega_v),
\]
其中 $\epsilon_v(s,\omega_v,\psi_v)$ 由命题 3.7 显式给出。注意几乎所有的 $\epsilon_v(s,\omega_v,\psi_v)$ 都是 $1$，故可定义\emph{整体 epsilon 因子}
\begin{equation*}\tag{4.6}
\epsilon(s,\omega)=\prod_v\epsilon_v(s,\omega_v,\psi_v),
\end{equation*}
并由局部者取乘积得到函数方程
\begin{equation*}\tag{4.7}
\widehat{z_o(1-s,\omega^{-1})}=\epsilon(s,\omega)\,z_o(s,\omega).
\end{equation*}
这里记号中略去了加法特征 $\psi$，因为事实上这个乘积已不再依赖所做的选择！这归功于各处的 (3.29) 式、每个整体加法特征对给定的非平凡 $\psi$ 都有形式 $\psi_\beta$（$\beta\in k^\times$）这一事实，以及乘积公式。于是
\begin{equation*}\tag{4.8}
\begin{aligned}
\Lambda(1-s,\omega^{-1})\cdot\widehat{z_o(1-s,\omega^{-1})}
&=\widehat{z(1-s,\omega^{-1})}\\
&=z(s,\omega)\\
&=\Lambda(s,\omega)\cdot z_o(s,\omega).
\end{aligned}
\end{equation*}
把关系 (4.7) 代入 (4.8) 的第一个表达式，得
\begin{equation*}\tag{4.9}
\Lambda(1-s,\omega^{-1})\cdot\epsilon(s,\omega)\cdot z_o(s,\omega)
=\Lambda(s,\omega)\cdot z_o(s,\omega).
\end{equation*}
但最后，由构造可知分布 $z_o(s,\omega)$ 处处不为零，于是得到\emph{函数方程}：

\medskip\noindent\textbf{推论 4.4.}\quad
\[
\Lambda(s,\omega)=\epsilon(s,\omega)\cdot\Lambda(1-s,\omega^{-1}).
\]

\begin{annote}
这一段的拼装顺序值得复述：定理 4.3 给出 $z$ 自身的函数方程，属于最难的解析输入（Poisson 求和）；(4.3) 把 $z$ 换成 $\Lambda z_o$；$z_o$ 整，故 $\Lambda$ 的延拓随之而来；局部 (3.23) 相乘得 (4.7)；代入 (4.8) 后 $z_o$ 可约去，才得推论 4.4。原文有两处交叉引用错误：称为 (3.13) 的局部函数方程实为 (3.23)，而「命题 3.7」应为命题 3.8（3.7 是推论）。
\par\vspace{0.15em}{\raggedleft 【来源：原书 p.~128（PDF 133）；Tate 论文 §4--§6（A）】\par}
\end{annote}

最后，定理 4.3 的证明归结为 Riemann 基于 Poisson 求和的经典论证！这一非常标准的计算，读者可参阅 Tate 论文 pp.~339--341、Weil [13, pp.~121--124]，或 Bump [2, pp.~267--270]。

\section{例子}

这里从简。

对数域 $k$ 与 $\A^\times/k^\times$ 的（酉）特征 $\omega=\otimes_v\omega_v$，令
\begin{equation*}\tag{5.1}
\mathfrak m=\prod_v\mathfrak p_v^{c_v(\omega)},
\end{equation*}
其中乘积取遍非阿基米德位，$c_v(\omega)=c_v(\omega_v)$ 是 $\omega_v$ 的导子。于是对有限位 $v$，$\mathfrak p_v\nmid\mathfrak m$ 当且仅当 $\omega$ 在 $v$ 处无分歧。我们可以定义 $I_{\mathfrak m}$（$k$ 的与 $\mathfrak m$ 互素的分数理想群）的特征 $\chi$：置
\[
\chi(\mathfrak p_v)=\omega_v(\varpi_v).
\]
由于 $\omega_v$ 无分歧，此值与局部生成元 $\varpi_v$ 的选取无关。令
\begin{equation*}\tag{5.2}
k^\times_{\mathfrak m}=\{\alpha\in k^\times\mid\alpha\equiv1\ \ (\mathrm{mod}\ \mathfrak m)\},
\end{equation*}
其中条件 $\alpha\equiv1\ (\mathrm{mod}\ \mathfrak m)$ 意即对一切满足 $c_v(\omega)>0$ 的 $v$ 有 $\operatorname{ord}_v(\alpha-1)\ge c_v(\omega)$，而在其它位对 $\alpha$ 不加条件。以 $P_{\mathfrak m}$ 记由 $\alpha\in k^\times_{\mathfrak m}$ 生成的（主分数理想）$(\alpha)$ 构成的群。由 $\omega(\alpha)=1$ 这一事实推出
\begin{equation*}\tag{5.3}
\chi((\alpha))=\omega_\infty(\alpha)^{-1},
\end{equation*}
其中 $(\alpha)$ 是由 $\alpha$ 生成的主理想，$\omega_\infty=\prod_{v\ \text{阿基米德}}\omega_v$ 是 $\omega$ 的阿基米德部分。于是特征 $\omega$ 决定 $I_{\mathfrak m}$ 的特征 $\chi$，$\chi$ 在 $P_{\mathfrak m}$ 上的限制由某个阿基米德特征在理想生成元上的取值给出。\footnote{我相信这正是 Hecke 用 Gr\"o\ss encharakter 一词所指的意思：即在 $(\alpha)\in P_{\mathfrak m}$ 上，$\chi$ 依赖阿基米德位 $v$ 处的量（Gr\"o\ss en）$\alpha_v\in k_v$。}由于 $P_{\mathfrak m}$ 在 $I_{\mathfrak m}$ 中指数有限，具有给定导子与给定阿基米德部分的这种特征只有有限多个。对固定的 $\mathfrak m$，这个阿基米德特征还受如下要求约束：若 $\epsilon\in k^\times_{\mathfrak m}\cap\cO_k^\times$ 是满足 $\epsilon\equiv1\ (\mathrm{mod}\ \mathfrak m)$ 的单位，则
\begin{equation*}\tag{5.4}
\omega_\infty(\epsilon)=1.
\end{equation*}

反之，任何对某个 $\mathfrak m$ 满足此条件的 $\omega_\infty$，都决定 $P_{\mathfrak m}$ 的一个特征，它可以延拓，给出 $I_{\mathfrak m}$ 的有限多个 Gr\"o\ss encharakter $\chi$。

当然，由经典 Hecke 特征 $\chi$ 可以复原 $\omega$。从这样一个具有给定阿基米德型 $\omega_\infty$（即满足 (5.3)）的 $\chi$ 出发，对 $\mathfrak p_v\nmid\mathfrak m$ 置 $\omega_v(\varpi_v)=\chi(\mathfrak p_v)$，定义 $k_v^\times$ 的无分歧特征 $\omega_v$。对 $\mathfrak p_v\mid\mathfrak m$ 的 $v$，$\omega$ 的其余分量 $\omega_v$ 由条件「对一切 $\alpha\in k^\times$ 有 $\omega(\alpha)=1$」唯一确定。更确切地说，弱逼近定理 [3] 表明：对任何有限位集 $S$，$k^\times$ 的像在乘积 $\prod_{v\in S}k_v^\times$ 中稠密。分量 $\omega_v$ 的导子至多为 $\operatorname{ord}_v(\mathfrak m)$；若 $\chi$ 模 $\mathfrak m$ 本原，则导子恰为 $\operatorname{ord}_v(\mathfrak m)$，如此等等。

\begin{annote}
这里是 Tate 论文与经典 Hecke 理论的接口，可小结为一张字典：$\mathfrak m$ 记导子，$I_{\mathfrak m}$ 是与之互素的分数理想群，$P_{\mathfrak m}$ 是「$\equiv1\bmod\mathfrak m$」的主理想子群，$I_{\mathfrak m}/P_{\mathfrak m}$ 就是射线类群。$\omega$ 的阿基米德部分 $\omega_\infty$ 在单位群上所受的约束 (5.4) 正是 $\chi$ 良定义的条件；反过来，给定 $\chi$ 后，非阿基米德分量由 $\chi(\mathfrak p_v)$ 定出，阿基米德分量由 (5.3) 定出，中间差的自由度正好被 (5.4) 与弱逼近消掉。原书在印 130 的练习 (iii) 末尾又用了一次编号 (5.4)，与这里的 (5.4) 重复，系原书编号疏漏。
\par\vspace{0.15em}{\raggedleft 【来源：Tate 论文 §3；Cassels--Fr\"ohlich 第 VII 章；Kowalski 本书第 2 章 [8]（A）】\par}
\end{annote}

对 $\omega_\infty=1$，可取任何 $\mathfrak m$。取 $\mathfrak m=\cO_k$ 时得到处处无分歧的特征，即理想类群 $H=I/P$ 的特征。对更一般的 $\mathfrak m$，得到的是 Kowalski 几章中提到的射线类群 $H_{\mathfrak m}=I_{\mathfrak m}/P_{\mathfrak m}$ 的特征。

$\omega=1$ 的情形当然给出域 $k$ 的 Dedekind zeta 函数。这个例子在 [13, 第 VII 章第 6 节] 中有详细解释，而更一般的 $\omega$ 的情形在该章第 7 节中解释。

\medskip\noindent\textbf{练习。}\quad 设 $k$ 是实二次域 $k=\mathbb Q(\sqrt5)$。则 $\cO_k=\mathbb Z+\mathbb Z\epsilon$，其中 $\epsilon=\frac{1+\sqrt5}{2}$ 是基本单位，且 $\cO_k^\times=\{\pm1\}\times\{\epsilon^n\mid n\in\mathbb Z\}$。

(i) 先设 $\mathfrak m=\cO_k$。则特征
\[
\omega_\infty(x_1,x_2)=\left|\frac{x_1}{x_2}\right|^{-i\pi/\log\epsilon}
\]
在 $\cO_k^\times$ 上平凡，从而由
\[
\chi((\alpha))=\omega_\infty(\alpha)^{-1}=\left|\frac{\alpha}{\alpha'}\right|^{i\pi/\log\epsilon}
\]
定义了 $I$ 的 Hecke 特征 $\chi$。写出相应的 Hecke $L$ 级数及其函数方程。

其次取 $\mathfrak m=(\sqrt5)=\mathfrak p_5$。

(ii) 验证：满足 $u\equiv1\ (\mathrm{mod}\ \mathfrak m)$ 的单位群 $k^\times_{\mathfrak m}\cap\cO_k^\times$ 由 $-\epsilon^2$ 生成。

(iii) 定出所有对 $\mathfrak m$ 满足 (5.3) 的特征 $\omega_\infty:\mathbb R^\times\times\mathbb R^\times\to\mathbb C^\times$。例如
\begin{equation*}\tag{5.4}
\omega_\infty(x_1,x_2)=\operatorname{sgn}(x_1)\left|\frac{x_1}{x_2}\right|^{-i\pi/(4\log\epsilon)}
\quad\text{与}\quad
\omega_\infty(x_1,x_2)=\left|\frac{x_1}{x_2}\right|^{-i\pi/(2\log\epsilon)}
\end{equation*}
就是两种可能。

(iv) 对给定的 $\omega_\infty$，定出 $\mathfrak p_5=\mathfrak m$ 处的局部成分 $\omega_5$。注意 $\cO_5^\times/(1+\mathcal P_5)\simeq\mathbb F_5^\times$，且整体单位映满这个群。又 $\varpi_5=\sqrt5$ 给出局部一致化子。例如，对由 (5.4) 给出的 $\omega_\infty$，$\omega_5$ 是分歧特征，满足 $\omega_5(\sqrt5)=1$ 且 $\omega_5(\epsilon)=i$。

(v) 定出与 $\omega$ 相配的 $I_{\mathfrak m}$ 的经典 Hecke 特征 $\chi$。

(vi) 写出相应的 Hecke $L$ 级数及其函数方程。

\subsection*{参考文献}

\begin{enumerate}[label={[\arabic*]},leftmargin=2.8em,itemsep=0.35em]
\item J. Bernstein and A. Zelevinski, \emph{Representations of the group $GL(n,F)$ where $F$ is a nonarchimedean local field}, Russian Math. Surveys, \textbf{3}(1976), 1--68.
\item D. Bump, \emph{Automorphic Forms and Representations}, Cambridge University Press, Cambridge, UK, 1997.
\item J. W. S. Cassels and A. Fr\"ohlich, eds., \emph{Algebraic Number Theory}, Thompson Book Company, Washington, DC, 1967.
\item F. G. Friedlander, \emph{Introduction to the Theory of Distributions}, Cambridge University Press, Cambridge, UK, 1982.
\item S. Gelbart, \emph{Automorphic Forms on Adele Groups}, Princeton University Press, Princeton, NJ, 1975.
\item S. Gelbart, An elementary introduction to the Langlands program, \emph{Bull. Amer. Math. Soc.}, \textbf{10}(1984), 177--219.
\item E. Hecke, Eine neue Art von Zetafunktionen und ihre Beziehungen zur Verteilung der Primzahlen, Erste und Zweite Mitteilung, in \emph{Mathematische Werke}, Vandenhoeck and Ruprecht, G\"ottingen, The Netherlands, 1959.
\item E. Kowalski, Elementary theory of $L$-functions I and II, in J. Bernstein and S. Gelbart, eds., \emph{An Introduction to the Langlands Program}, Birkh\"auser Boston, Boston, 2003, 1--37（本书）.
\item S. Lang, \emph{Algebraic Number Theory}, Addison--Wesley, Reading, MA, 1970.
\item J. Tate, \emph{Fourier Analysis in Number Fields and Hecke's Zeta-Functions}, Ph.D. thesis, Princeton University, Princeton, NJ, 1950.
\item R. Valenza and D. Ramakrishnan, \emph{Fourier Analysis on Number Fields}, Graduate Texts in Mathematics 186, Springer-Verlag, New York, 1999.
\item A. Weil, Fonction zeta et distributions, \emph{S\'eminar Bourbaki}, \textbf{III}(1966); also appears in [14].
\item A. Weil, \emph{Basic Number Theory}, Springer-Verlag, New York, 1967.
\item A. Weil, \emph{Collected Papers}, Springer-Verlag, New York, 1979.
\end{enumerate}
